% SGA 3，第 VIII 篇，英文重建底本说明。
% 范围：第 7 节开头至推论 7.4。
% 权威源：Polo--Gille Exp8-8nov09.pdf, 组件内第 23--26 页,
% 完整合阅本第 639--642 页.
% 英文底本说明：当时的 Reinhold 英译本仅作对照材料；其自身版本
% 声明指出第 7 节需要修订。英文底本声称，此处每一公式和陈述均已
% 与 Polo--Gille 权威源核对。此为英文底本的历史说明。

\setcounter{footnote}{31}
\SGARefTarget{sga3:viii:section:7}\section*{7. 附录：关于群预概形的单态射}
\addcontentsline{toc}{section}{7. 附录：关于群预概形的单态射}

% 源页码：Exp8 组件内第 23 页; 完整合阅本第 639 页.
除 X，\SGARefLink{sga3:x:theorem:8-8}{8.8} 及第 XV、XVI 讲外，本讨论班其余部分都不需要\SGARefLink{sga3:viii:section:7}{本节}所证的结果。它实质上依赖于第 VI\(_A\) 讲中商群的存在定理。
\footnote{编者注：参见第 VI\(_A\) 讲，\SGARefLink{thm:sga3-VIA-3.2}{定理 3.2} 和 \SGARefLink{thm:sga3-VIA-3.3.2}{3.3.2}。}

\SGARefLink{sga3:viii:corollary:5:7}{推论 5.7} 引出了如下问题：何时可以断言 \(S\)-群的单态射 \(u:G\to H\) 为浸入，甚至为闭浸入。

在 VI\(_B\)，\SGARefLink{cor:sga3-VIB-1.4.2}{1.4.2} 中已见，当 \(S\) 为域 \(k\) 的谱时，只要 \(G\) 在 \(k\) 上有限型，且 \(H\) 在 \(k\) 上局部有限型，就有此结论。容易推出，仅假设 \(S\) 为 Artin 时，同一结果仍成立。\footnote{编者注：应考虑到第 VI\(_B\) 讲中所作的增补。}

另一方面，很容易给出并非浸入的双射单态射的例子，例如在域上的仿射直线或离散赋值环的谱上。取
\[
  H=(\mathbf Z/2\mathbf Z)_S,\qquad G=G_1\times_SG_2,
\]
其中 \(G_1\) 为 \(H\) 中除去纤维 \(H_s\simeq\mathbf Z/2\mathbf Z\) 中的闭点 \(x\ne0\) 所得的开子群，\(s\) 为 \(S\) 的一个固定闭点；\(G_2\) 为 \(H\) 的闭子概形，由单位截面给出的子概形与 \(x\) 定义的约化闭子概形之并构成。立即可验证，\(G_2\) 在 \(H\) 的乘法下稳定，因而为群概形。浸入 \(G_i\to H\)、\(i=1,2\) 便定义了 \(S\)-群同态
\[
  G=G_1\times_SG_2\longrightarrow H
\]
它显然是双射单态射，甚至为局部浸入，但不是浸入。事实上，\(G\) 与 \(H\) 都约化；\(G\) 有三个不交的不可约分支，而 \(H\) 只有两个。

还应注意，\(G_1\to H\) 给出了 \(H\) 的一个不闭的开子群，这与域上的代数群的情形不同。椭圆曲线退化理论还给出了此现象的其他例子，其中 \(H\) 在 \(S\) 上平滑，相对维数为一，且 \(G\) 的纤维连通。

人们或许会希望\begingroup
\renewcommand{\thefootnote}{(\(\ast\))}%
\footnote[1]{事实上，M. Raynaud 构造了一个 \(G\) 平滑且纤维连通的反例；参见 XVI，\SGARefLink{sga3:xvi:subsection:1:1:item:c}{1.1(c)}。若无纤维连通的假设，可取 \(S=\operatorname{Spec}(\mathbf Z_2)\)，令 \(G=(\mathbf Z/2\mathbf Z)_S\) 除去其非单位闭点，并令 \(H=(\mu_2)_S\)。}%
\endgroup
一旦假设 \(G\) 在 \(S\) 上平坦，且例如 \(G\) 和 \(H\) 在 \(S\) 上有限呈示，则每个 \(S\)-群的单态射 \(u:G\to H\) 都自动为浸入。我们将在附加假设下证明一个这种类型的结果。

首先，可以假设 \(S\) 仿射，并利用 \(G\) 和 \(H\) 的有限呈示假设，假设 \(S\) 为 Noether：归结到 \(\mathbf Z\) 上有限型环的谱。\footnote{编者注：参见 EGA~IV\(_3\)，§8 和第 VI\(_B\) 讲，§10。}于是 \(G\) 和 \(H\) 都是 Noether 的。此时，说 \(u:G\to H\) 为闭浸入（分别为浸入），等价于说 \(u\) 为单态射——这由假设成立——并且固有（分别在 \(u(G)\) 的每一点处固有）。\footnote{编者注：参见 EGA~IV\(_3\)，8.11.5 和 15.7.1。}

% 源页码：Exp8 组件内第 24 页; 完整合阅本第 640 页.
固有性的赋值判据表明，只须在每次基变换 \(S'\to S\) 后验证态射 \(u':G'\to H'\) 具有相应的固有性，其中 \(S'\) 为离散赋值环的谱，如愿还可要求完备。局部固有性的情形在 EGA~II，7.3 中遗漏，原拟作为勘误收入 EGA~IV。\begingroup
\renewcommand{\thefootnote}{(\(\ast\))}%
\footnote[0]{参见 EGA~IV\(_3\)，15.7。}%
\endgroup
因此，只要对 \(G,H,u\) 所加的任何附加假设在基变换下稳定，就归结到 \(S\) 本身为完备离散赋值环的谱的情形。

分别用 \(s\) 和 \(s'\) 表示 \(S\) 的闭点和泛点。纤维上诱导的同态
\[
  u_s:G_s\longrightarrow H_s,
  \qquad
  u_{s'}:G_{s'}\longrightarrow H_{s'},
\]
是闭浸入，因为它们是域上有限型群概形的单态射（VI\(_B\)，\SGARefLink{cor:sga3-VIB-1.4.2}{1.4.2}）。所以可将 \(G_{s'}\) 与 \(H_{s'}\) 的闭子概形等同。我们将用到\SGARefLink{sga3:viii:lemma:7:1}{如下结果}。

\medskip
\noindent\SGARefTarget{sga3:viii:lemma:7:1}\textbf{引理 7.1.}
\textit{设 \(S\) 为离散赋值环的谱，泛点为 \(s'\)。设 \(H\) 为 \(S\)-预概形，\(L'\) 为泛纤维 \(H_{s'}\) 的闭子预概形，因而也是 \(H\) 的子预概形。它在 \(H\) 中的概形论闭包 \(L\)——即 \(H\) 中包含 \(L'\) 的最小闭子预概形，参见 EGA~I，9.5——存在。它还是 \(H\) 中在 \(S\) 上平坦且泛纤维为 \(L'\) 的唯一闭子预概形。\(L\) 的构造关于对 \((H,L')\) 是函子性的，且与 \(S\) 上的纤维积交换。特别地，若 \(H\) 为 \(S\)-群，\(L'\) 为 \(H_{s'}\) 的子群，则 \(L\) 为 \(H\) 的子群。}

证明直接，留给读者。\begingroup
\renewcommand{\thefootnote}{(\(\ast\ast\))}%
\footnote[0]{参见 EGA~IV\(_2\)，2.8。}%
\endgroup
将此用于对 \((H,G_{s'})\)。单态射 \(u:G\to H\) 分解为
\[
  G\xrightarrow{\,u'\,}L\longrightarrow H,
\]
其中 \(L\) 为 \(H\) 的子群，是在 \(S\) 上平坦的闭子预概形，且 \(u':G\to L\) 在泛纤维上诱导同构。因此，\(u:G\to H\) 为浸入（分别为闭浸入），当且仅当 \(u':G\to L\) 也是如此。于是归结到 \(H\) 在 \(S\) 上平坦且 \(u_{s'}:G_{s'}\to H_{s'}\) 为同构的情形；仍要求将 \(H\) 换成闭子群预概形后，附加假设保持成立。

由于此时 \(H\) 是 \(H_{s'}\) 的概形论闭包，若 \(u\) 为浸入，其像 \(u(G)\) 为 \(H\) 中包含 \(H_{s'}\) 的子预概形。\footnote{编者注：删去了“闭”一词。}其概形论闭包因而为 \(H\)，所以 \(u(G)\) 为 \(H\) 的开子预概形。因此，必须证明 \(u\) 为开浸入；若所求结论为闭浸入，则必须证明它为同构。

% 源页码：Exp8 组件内第 25 页; 完整合阅本第 641 页.
由于 \(G\) 和 \(H\) 在 \(S\) 上平坦，这等价于纤维上诱导的映射为开浸入、分别为同构（参见 SGA~1，I，\SGARefExternal{sga3:viii:corollary:5:7}{5.7}）。这对 \(u_{s'}\) 已成立，所以只须证明 \(u_s\) 为开浸入、分别为同构。

由于 \(G\) 和 \(H\) 在 \(S\) 上平坦，其纤维维数恒定（VI\(_B\)，\SGARefLink{cor:sga3-VIB-4.3}{4.3}）。它们的泛纤维相同，因此 \(G\) 和 \(H\) 的这些维数相等。于是 \(u_s:G_s\to H_s\) 为相同维数的代数群的单态射。由此可知，\(G_s\) 在 \(H_s\) 中开，且在集合意义下为 \(H_s\) 的若干连通分支之并。事实上，转到代数闭基域并将 \(G\) 和 \(H\) 约化后，二者在该域上都平滑，断言随即成立。

补集 \(H_s-G_s\) 在 \(H_s\) 中闭，因而在 \(H\) 中闭。它在 \(H\) 中的补集 \(H'\) 是开的，且在群运算下稳定。因此，用 \(H'\) 代替 \(H\)，在通常的保留条件下可归结到 \(u\) 为双射，即 \(u_s:G_s\to H_s\) 为双射的情形。无论如何，都归结到证明 \(u\)，等价地 \(u_s\)，为同构，可能还带有这一双射性假设。

先假设 \(u\) 为双射。若 \(H_s\) 约化，则 \(u_s\) 为同构，因为 \(G_s\) 与 \(H_s\) 中具有相同底集合的闭子概形等同。特别地，若 \(k=\kappa(s)\) 的特征为零，则由 Cartier 定理，\(k\) 上的每个代数群都约化（参见 VI\(_B\)，1.6.1；VII\(_B\)，3.3.1；或 EGA~IV\(_4\)，16.12.2 和 17.12.5）。我们证明了：

\medskip
\noindent\SGARefTarget{sga3:viii:proposition:7:2}\textbf{命题 7.2.}
\textit{设 \(u:G\to H\) 为 \(S\) 上有限呈示群预概形的同态。若 \(u\) 为单态射，\(G\) 在 \(S\) 上平坦，且 \(S\) 的每个剩余域的特征均为零，则 \(u\) 为浸入。}

现在假设 \(\kappa(s)\) 的特征为 \(p>0\)，并限制在 \(G\) 交换的情形。经过上述归约，\(H\) 也交换，因为它是 \(H_{s'}\simeq G_{s'}\) 的概形论闭包。对 \(n\geq0\)，置
\[
  S_n=\operatorname{Spec}(V/\mathfrak m^{n+1}),\qquad
  G_n=G\times_SS_n,\qquad H_n=H\times_SS_n,
\]
其中 \(V\) 为定义 \(S\) 的赋值环，\(\mathfrak m\) 为其极大理想。对 \(m>0\)，用 \({}_mG\) 和 \({}_mH\) 表示 \(G\) 和 \(H\) 中乘以 \(m\) 的核。同样定义
\[
  {}_m(G_n)=({}_mG)_n,
\]
简记为 \({}_mG_n\)，并对 \(H\) 作同样约定。

由 VI\(_A\)，\SGARefLink{thm:sga3-VIA-3.2}{3.2}，商
\[
  Q_n=H_n/G_n
\]
存在。每个 \(Q_n\) 都是 \(S_n\) 上的交换群概形，在 \(S_n\) 上平坦，且 \(H_n\to Q_n\) 忠实平坦，核为 \(G_n\)。商的构造与基扩张交换，\footnote{编者注：应在第 V 和 VI\(_A\) 讲中明确指出这一点。}因此
\[
  Q_n\simeq Q_m\times_{S_m}S_n\qquad(m\geq n).
\]
特别地，\(Q_n\times_{S_n}S_0=Q_0=H_0/G_0\)。由于 \(G_0\) 和 \(H_0\) 有相同的底集合，\(Q_0\) 在集合意义下为一个点：它是纯无穷小群。

% 源页码：Exp8 组件内第 26 页; 完整合阅本第 642 页.
由此可知，每个 \(Q_n\) 都在 \(S_n\) 上有限且平坦。设其由 \(V_n=V/\mathfrak m^{n+1}\) 上的代数 \(C_n\) 定义，后者作为 \(V_n\)-模有限自由。对 \(m\geq n\)，有
\[
  C_n=C_m\otimes_{V_m}V_n,
\]
且与余乘法映射相容。因此
\[
  C=\varprojlim C_n
  \quad\text{为下列环上的有限自由模}\quad
  V=\varprojlim V_n,
\]
各 \(C_n\) 上的余乘法在 \(C\) 上诱导余乘法。于是 \(Q=\operatorname{Spec}(C)\) 为 \(S\) 上的有限平坦群概形，对每个 \(n\) 满足
\[
  Q\times_SS_n=Q_n
\]

\medskip
\noindent\SGARefTarget{sga3:viii:lemma:7:3}\textbf{引理 7.3.}
\footnote{编者注：这里删去了 \(Q\) 交换的假设，并相应修改了 \SGARefLink{sga3:viii:remark:7:3:1}{7.3.1}。若 \(Q\) 不交换，\(n\) 次幂态射未必为群同态。}
\textit{设 \(K\) 为域，\(Q\) 为 \(K\) 上次数为 \(n\) 的有限群概形。则 \(Q\) 上的 \(n\) 次幂态射为零。}

见 VII\(_A\)，\SGARefLink{sga3:VIIA:8.5}{8.5}。

\medskip
\noindent\SGARefTarget{sga3:viii:remark:7:3:1}\textbf{注 7.3.1.}
引理~\SGARefLink{sga3:viii:lemma:7:3}{7.3} 的陈述对任意基预概形 \(S\) 上的有限局部自由群概形 \(Q/S\) 仍有意义。若能在此一般性下找到证明，将很有趣。

引理~\SGARefLink{sga3:viii:lemma:7:3}{7.3}（即 VII\(_A\)，8.5）确实表明，当 \(S\) 约化时，此陈述成立：将它用于 \(Q\) 在 \(S\) 的极大点，即其不可约分支的泛点处的纤维即可。\footnote{编者注：应将此补入第 VII\(_A\) 讲。另一方面，当 \(Q\) 交换时，该陈述对每个 \(S\) 都成立；参见 \SGARefLink{sga3:viii:bib:to70}{[TO70]}，第~4~页的 Deligne 定理。}

在前面的情形中，\(S\) 为离散赋值环的谱，\(Q\) 交换，于是得到 \(n\operatorname{id}_Q=0\)。这里 \(n\) 为剩余特征的一个幂 \(p^{\nu_0}\)。因此：

\medskip
\noindent\SGARefTarget{sga3:viii:corollary:7:4}\textbf{推论 7.4.}
\textit{设 \(Q\) 和各 \(Q_n\) 如上，其共同次数为 \(p^{\nu_0}\)。则}
\[
  p^{\nu_0}\operatorname{id}_Q=0,
  \qquad
  p^{\nu_0}\operatorname{id}_{Q_n}=0
  \quad\textit{对每个 \(n\)}.
\]
